% ===========================================================================
% Chapter 3: FusionCropNet V6 Architecture
% ===========================================================================
\section{FusionCropNet V6 Architecture}

\subsection{Overview}

FusionCropNet V6 is a hierarchical multi-scale multi-task architecture that
implements all four mathematical theory modules through 14 components. It
inherits the EDL (Evidential Deep Learning) framework from V5EDL and activates
all V6 enhancements by default.

\begin{figure}[h]
\centering
\begin{tikzpicture}[node distance=0.8cm, auto,
    block/.style={rectangle, draw, minimum width=2.5cm, minimum height=0.7cm,
    font=\footnotesize, rounded corners=2pt},
    math/.style={block, fill=blue!10},
    arrow/.style={->, >=stealth, thick}]

    % Input
    \node[block] (input) {INPUT: opt(B,T,10,H,W) sar(B,T,5,H,W) dem(B,5,H,W)};

    % Early fusion
    \node[block, below=of input] (early) {1. Early Fusion: ModalNorm + Concat(20ch)};

    % Encoders
    \node[block, below=of early] (enc) {2. Encoders: Optical(SeCo) + SAR(DEM-FiLM) + DEM};

    % DEM 5 paths
    \node[block, below=of enc] (dem) {3. DEM 5-Path Injection};

    % Temporal
    \node[block, below=of dem] (temp) {4. Multi-Scale Temporal: TemporalLite $\times$ 3};

    % Cross-modal
    \node[block, below=of temp] (cross) {5. 3-Scale Cross-Modal Attention};

    % Late fusion
    \node[block, below=of cross] (late) {6. LateFusion + Decoder (CARAFE)};

    % Heads
    \node[block, below=of late] (heads) {7. Multi-Task Heads: EDL + LAI + Growth + Boundary};

    % Math modules (connected vertically on right side)
    \node[math, right=3cm of enc] (m1) {Module 1: Diff. Geometry};
    \node[math, right=3cm of dem] (m2) {Module 2: Optimal Transport};
    \node[math, right=3cm of cross] (m3) {Module 3: Alg. Topology};
    \node[math, right=3cm of temp] (m4) {Module 4: Stat. Learning};

    \draw[arrow] (input) -- (early);
    \draw[arrow] (early) -- (enc);
    \draw[arrow] (enc) -- (dem);
    \draw[arrow] (dem) -- (temp);
    \draw[arrow] (temp) -- (cross);
    \draw[arrow] (cross) -- (late);
    \draw[arrow] (late) -- (heads);

    \draw[arrow, dashed, red] (m1) -- (dem);
    \draw[arrow, dashed, red] (m2) -- (cross);
    \draw[arrow, dashed, red] (m3) -- (cross);
    \draw[arrow, dashed, red] (m4) -- (temp);
\end{tikzpicture}
\caption{FusionCropNet V6 architecture with mathematical theory module mapping.}
\end{figure}

\subsection{Component Details}

\subsubsection{DEM Encoder (Module 1 Integration)}

The standard CNN DEM encoder is replaced with explicit geometric invariant
extraction:
\begin{equation}
    \text{DEM}(5\text{ch}) \xrightarrow{\text{SIREN fit}} h(x,y)
    \xrightarrow{\text{Eqs. \eqref{eq:K}--\eqref{eq:tau}}}
    \{K, H, \kappa_1, \kappa_2, \tau_g\}(5\text{ch})
    \xrightarrow{\text{light CNN}} \mathbf{f}_{\text{dem}}(128\text{ch})
\end{equation}

The geometric semantics of each channel:
\begin{itemize}
    \item $K$: elliptic ($K>0$, peaks/valleys) vs. hyperbolic ($K<0$, saddles)
    vs. parabolic ($K\approx0$, slopes)
    \item $H$: local mean bending (complements slope)
    \item $\kappa_1, \kappa_2$: asymmetric slope characterization
    \item $\tau_g$: surface torsion (critical for terrace/terrace identification)
\end{itemize}

\subsubsection{Five DEM Injection Paths}

DEM modulates all processing stages through five paths:

\begin{enumerate}
    \item \textbf{DEM $\to$ SAR FiLM:} $\mathbf{s}_i' = \gamma_i(\mathbf{d}) \odot \mathbf{s}_i + \beta_i(\mathbf{d})$.
    Curvature-driven scatter correction.

    \item \textbf{DEM $\to$ Optical FiLM:} Terrain shadow and atmospheric
    thickness correction via curvature-dependent reflectance model.

    \item \textbf{DEM $\to$ Temporal FiLM:} Elevation + curvature joint
    phenological offset: $\Delta T \approx -0.6^\circ\text{C}/100\text{m}$.

    \item \textbf{DEM $\to$ Decoder Skip:} Explicit geometric prior at full
    resolution via $1\times1$ convolution.

    \item \textbf{DEM $\to$ Spatial Conditioner:} Terrain-aware gating of
    cross-modal fusion: strengthen in peaks ($K>0$), weaken in saddles ($K<0$).
\end{enumerate}

\subsubsection{TemporalLite (Module 4 Integration)}

Three TemporalLite instances encode multi-scale temporal sequences:
\begin{equation}
    \text{TemporalLite}(\mathbf{x}) = \gamma \cdot \frac{1}{T}\sum_{t=1}^T
    \text{LayerNorm}(\text{DepthwiseConv1D}(\mathbf{x}, \mathbf{W}_c))
\end{equation}

Spectral-Balanced Initialization (Theorem 3, Module 4) ensures gradient
stability without warmup. Complexity: $O(T \cdot D)$ vs. Transformer
$O(T \cdot D^2 + T^2 \cdot D)$. For $D=64$: only 321 parameters.

\subsubsection{Cross-Modal Attention (Module 2 Integration)}

Three-scale cross-modal attention with Grassmann subspace alignment:
\begin{equation}
    \mathbf{f}_H = \text{CrossModalLite}(\mathbf{o}_{\text{aligned}}, \mathbf{s}_{\text{aligned}})
\end{equation}
where features are pre-aligned via joint OT+Grassmann optimization
(Eq.~\eqref{eq:joint-ot}) before QKV projection.

\subsubsection{EDL Head (Module 3 Integration)}

The EDL head outputs Dirichlet parameters $\boldsymbol{\alpha} = \text{softplus}(\mathbf{z}) + 1$.
The vacuity $u = K / \sum_k \alpha_k$ is replaced with topological conflict
classification:
\begin{equation}
    \text{conflict\_type} = \begin{cases}
        \text{Noise} & [\omega_1 \smile \omega_2] = 0, \kappa < \tau_\kappa \\
        \text{Structural} & [\omega_1 \smile \omega_2] \neq 0 \\
        \text{HighOrder} & \exists p \geq 1, H^p(K) \ni [\eta] \neq 0
    \end{cases}
\end{equation}

\subsubsection{Multi-Task Pseudo-Label Learning}

Four auxiliary tasks provide self-supervised signals:
\begin{itemize}
    \item \textbf{LAI Regression:} $LAI \approx -\ln((\text{NDVI}_{\max} - \text{NDVI})/
    (\text{NDVI}_{\max} - \text{NDVI}_{\min})) / k$, Huber loss
    \item \textbf{Growth Stage:} DOY discretization into 5 phenological stages
    \item \textbf{Boundary Detection:} DEM Sobel gradient thresholding, BCE + Dice loss
    \item \textbf{Scene Understanding:} LightSceneHead (18K params) predicts
    paddy/dryland/mixed/facility types
\end{itemize}

Multi-task loss uses homoscedastic uncertainty weighting (Kendall et al. 2018):
\begin{equation}
    \mathcal{L}_{\text{total}} = \sum_i \frac{1}{\sigma_i^2}\mathcal{L}_i + \log\sigma_i
\end{equation}

\subsection{Forward Propagation Algorithm}

\begin{algorithm}[h]
\caption{FusionCropNet V6 Forward Propagation}
\label{alg:forward}
\begin{algorithmic}[1]
\Require $\mathbf{X}_{\text{opt}} \in \mathbb{R}^{B \times T \times 10 \times H \times W}$,
         $\mathbf{X}_{\text{sar}} \in \mathbb{R}^{B \times T \times 5 \times H \times W}$,
         $\mathbf{X}_{\text{dem}} \in \mathbb{R}^{B \times 5 \times H \times W}$,
         $\mathbf{doy} \in [0,1]^{B \times T}$
\Ensure $\boldsymbol{\alpha} \in \mathbb{R}^{B \times K \times H \times W}$,
        $\{y_{\text{LAI}}, y_{\text{growth}}, y_{\text{boundary}}\}$

\State \textbf{// Step 1: Geometric Invariant Extraction (Module 1)}
\State $\mathbf{g} \gets \text{compute\_geometric\_invariants}(\mathbf{X}_{\text{dem}})$
    \Comment{5 ch: $\{K, H, \kappa_1, \kappa_2, \tau_g\}$}
\State $\mathbf{f}_{\text{dem}} \gets \text{GeometricInvariantEncoder}(\mathbf{X}_{\text{dem}})$
    \Comment{128 ch}

\State \textbf{// Step 2: Early Fusion}
\State $\mathbf{u}_{\text{early}} \gets \text{Concat}(\text{LN}(\mathbf{X}_{\text{opt}}[0]),
\text{LN}(\mathbf{X}_{\text{sar}}[0]), \text{LN}(\mathbf{X}_{\text{dem}}))$
    \Comment{20 ch}
\State $\mathbf{u}_{\text{proj}} \gets \text{Conv}_{1\times1}^{20 \to 128}(\mathbf{u}_{\text{early}})$

\State \textbf{// Step 3: Encoder Forward Pass}
\State $\mathbf{f}_{\text{opt}}, \mathbf{o}_{\text{p2}}, \mathbf{o}_{\text{p3}} \gets
\text{OpticalEncoder}(\mathbf{X}_{\text{opt}})$ \Comment{SeCo backbone + FPN}
\State $\mathbf{s}_1, \mathbf{s}_2, \mathbf{s}_3 \gets
\text{SAREncoder}(\mathbf{X}_{\text{sar}}, \mathbf{f}_{\text{dem}})$
    \Comment{DEM FiLM at s1, s2}

\State \textbf{// Step 4: DEM Multi-Path Injection}
\State $\mathbf{f}_{\text{opt}}, \mathbf{o}_{\text{p2}} \gets
\text{DEMOpticalConditioner}(\mathbf{f}_{\text{opt}}, \mathbf{o}_{\text{p2}}, \mathbf{f}_{\text{dem}})$
\State $\mathbf{s}_1 \gets \text{FiLM}_{\text{sar}}(\mathbf{s}_1, \mathbf{f}_{\text{dem}})$
\State $\mathbf{s}_2 \gets \text{FiLM}_{\text{sar}}(\mathbf{s}_2,
\text{Interp}(\mathbf{f}_{\text{dem}}, \mathbf{s}_2))$
\State $\mathbf{b}_{\text{temp}} \gets \text{MLP}(\text{AvgPool}(\mathbf{f}_{\text{dem}}))$
    \Comment{Temporal elevation bias}

\State \textbf{// Step 5: Multi-Scale Temporal Encoding (Module 4)}
\State $\mathbf{o}_{\text{p2}}^t \gets \text{TemporalLite}_{256}(\mathbf{o}_{\text{p2}})$
    \Comment{Spectral-Balanced init}
\State $\mathbf{s}_1^t \gets \text{TemporalLite}_{64}(\mathbf{s}_1)$
\State $\mathbf{s}_2^t \gets \text{TemporalLite}_{128}(\mathbf{s}_2)$
\State $\mathbf{f}_{\text{opt}}^g, \mathbf{s}_3^g \gets
\text{Transformer}(\mathbf{f}_{\text{opt}}, \mathbf{s}_3)$
    \Comment{H/4 scale, gradient ckpt}

\State \textbf{// Step 6: Cross-Modal Alignment (Module 2)}
\State $\mathbf{f}_{\text{opt}}^a, \mathbf{s}_3^a \gets
\text{JointOTGrassmannAligner}(\mathbf{f}_{\text{opt}}^g, \mathbf{s}_3^g)$
    \Comment{SGW + Stiefel ADMM, $\sim$0.42s}

\State \textbf{// Step 7: Multi-Scale Cross-Attention}
\State $\mathbf{f}_H \gets \text{CrossModalLite}_{64}(\mathbf{o}_{\text{p2}}^t, \mathbf{s}_1^t)$
\State $\mathbf{f}_{H/2} \gets \text{CrossModalLite}_{128}(\mathbf{o}_{\text{p2}}^t, \mathbf{s}_2^t)$
\State $\mathbf{f}_{H/4} \gets \text{CrossModalAttention}_{512}(\mathbf{f}_{\text{opt}}^a,
\mathbf{s}_3^a)$

\State \textbf{// Step 8: Hierarchical Merge + Decoder}
\State $\mathbf{f}_{H/2} \gets \text{ConvMerge}(\mathbf{s}_2^t,
\mathbf{f}_H^{\downarrow 2\times})$
\State $\mathbf{f}_{H/4} \gets \text{ConvMerge}(\mathbf{s}_3^a,
\mathbf{f}_{H/2}^{\downarrow 2\times})$
\State $\mathbf{f}_{\text{pre}} \gets \text{Decoder}(\mathbf{f}_{H/4},
\{\mathbf{o}_{\text{p2}}^t\}, \{\mathbf{f}_H, \mathbf{f}_{H/2}\},
\mathbf{f}_{\text{dem}}, \mathbf{u}_{\text{proj}})$
    \Comment{CARAFE up + SpatialRefinement + DEM skip + Early skip}

\State \textbf{// Step 9: EDL Classification Head (Module 3)}
\State $\mathbf{z} \gets \text{EDLHead}(\mathbf{f}_{\text{pre}})$
\State $\boldsymbol{\alpha} \gets \text{softplus}(\mathbf{z}) + 1$
    \Comment{Dirichlet parameters, $(B, K, H, W)$}

\State \textbf{// Step 10: Topological Conflict Detection (Module 3)}
\State $(\kappa, [\omega]) \gets \text{detect\_conflict}(\boldsymbol{\alpha}_{\text{opt}},
\boldsymbol{\alpha}_{\text{sar}}, \boldsymbol{\alpha}_{\text{fused}})$
\State \textbf{if} $[\omega] \neq 0$ \textbf{then}
    $\boldsymbol{\alpha} \gets \text{persistent\_correction}(\boldsymbol{\alpha})$

\State \textbf{// Step 11: Multi-Task Auxiliary Heads}
\State $y_{\text{LAI}} \gets \text{LAIHead}(\mathbf{f}_{\text{pre}})$
\State $y_{\text{growth}} \gets \text{GrowthStageHead}(\mathbf{f}_{\text{pre}})$
\State $y_{\text{boundary}} \gets \text{BoundaryHead}(\mathbf{f}_{\text{pre}})$
\State $y_{\text{scene}} \gets \text{LightSceneHead}(\mathbf{f}_{\text{pre}})$

\State \Return $\boldsymbol{\alpha}, (y_{\text{LAI}}, y_{\text{growth}},
y_{\text{boundary}}, y_{\text{scene}})$
\end{algorithmic}
\end{algorithm}

\subsection{Unified Loss Function}

The complete training objective combines all four mathematical modules:

\begin{equation}\label{eq:unified-loss}
    \mathcal{L}_{\text{V6}} = \underbrace{\mathcal{L}_{\text{EDL}}}_{\text{Module 3}} +
    \lambda_1 \underbrace{\mathcal{L}_{\text{NDVI}}}_{\text{auxiliary}} +
    \lambda_2 \underbrace{\mathcal{L}_{\text{LAI}}}_{\text{phenology}} +
    \lambda_3 \underbrace{\mathcal{L}_{\text{growth}}}_{\text{phenology}} +
    \lambda_4 \underbrace{\mathcal{L}_{\text{boundary}}}_{\text{spatial}} +
    \lambda_5 \underbrace{\mathcal{L}_{\text{consistency}}}_{\text{Module 4}} +
    \lambda_6 \underbrace{\mathcal{L}_{\text{Topo-EWC}}}_{\text{Module 1 (TTA)}}
\end{equation}

where each component is:

\subsubsection{EDL Loss (Module 3)}

\begin{equation}
    \mathcal{L}_{\text{EDL}}(\boldsymbol{\alpha}, y) =
    \underbrace{\sum_{k=1}^K y_k(\psi(S) - \psi(\alpha_k))}_{\text{Type-II ML}} +
    \lambda_t \underbrace{\text{KL}(\text{Dir}(\tilde{\boldsymbol{\alpha}}) \|
    \text{Dir}(\mathbf{1}))}_{\text{KL regularization}}
\end{equation}
with $\lambda_t = \lambda_{\max} \cdot \min(1, t / T_{\text{anneal}})$.
The KL term encourages vacuity for misclassified pixels, calibrating
uncertainty estimates through Dirichlet shrinkage toward the uniform prior.

\subsubsection{Temporal Consistency Loss (Module 4)}

\begin{equation}
    \mathcal{L}_{\text{consistency}} = \frac{1}{|\mathcal{V}|}\sum_{i \in \mathcal{V}}
    \|\text{sigmoid}(\mathbf{w}_c^\top \mathbf{h}_i) -
      \text{sigmoid}(\bar{\mathbf{w}}_c^\top \bar{\mathbf{h}}_i)\|^2
\end{equation}
where $\mathcal{V}$ is the set of unmasked (non-cloud) temporal observations,
and $\bar{\cdot}$ denotes the target (EMA) projection.

\subsubsection{Multi-Task Weighting}

Following Kendall et al.~\cite{kendall2018multi}, task weights are learned
via homoscedastic uncertainty:
\begin{equation}
    \mathcal{L}_{\text{total}} = \sum_{i} \frac{1}{2\sigma_i^2}\mathcal{L}_i
    + \log\sigma_i \quad (\text{regression tasks})
\end{equation}
\begin{equation}
    \mathcal{L}_{\text{total}} = \sum_{i} \frac{1}{\sigma_i^2}\mathcal{L}_i
    + \log\sigma_i \quad (\text{classification tasks})
\end{equation}
where $\sigma_i = \exp(\log\_\text{var}_i)$ is a learned per-task uncertainty
parameter. This automatically balances the five tasks without manual weight tuning.

\subsubsection{Topo-EWC Regularization (Module 1 Level 2)}

During TTA, the loss is augmented with topological elastic weight consolidation:
\begin{equation}
    \mathcal{L}_{\text{Topo-EWC}} = \frac{\lambda_{\text{ewc}}}{2}
    \sum_i \Omega_i (\theta_i - \theta_{0,i})^2
\end{equation}
where $\Omega_i = \sum_{\sigma \in K} w_\sigma \cdot |\partial\sigma/\partial\theta_i|^2$
is the parameter importance derived from persistent homology barcode weights
$w_\sigma = |b_\sigma - d_\sigma|$.

\subsection{Training Protocol}

\begin{algorithm}[h]
\caption{Two-Phase Training Protocol}
\label{alg:training}
\begin{algorithmic}[1]
\Require Dataset $\mathcal{D}$, model $f_\theta$, max epochs $E_1, E_2$

\State \textbf{// Phase 1: Encoder Warm-up (frozen backbone)}
\For{epoch $e = 1, \ldots, E_1$}
    \For{batch $\mathbf{b} \in \mathcal{D}_{\text{train}}$}
        \State Freeze OpticalEncoder backbone, SAREncoder stem
        \State $\boldsymbol{\alpha}, \mathbf{y}_{\text{aux}} \gets
        f_\theta(\mathbf{b})$ \Comment{Algorithm~\ref{alg:forward}}
        \State $\mathcal{L} \gets$ Eq.~\eqref{eq:unified-loss}
        \State $\theta \gets \theta - \eta \nabla_\theta\mathcal{L}$
        \State \textbf{// No warmup needed: Spectral-Balanced Init (Theorem 3, Module 4)}
    \EndFor
\EndFor

\State \textbf{// Phase 2: Full Fine-tuning}
\For{epoch $e = E_1+1, \ldots, E_1+E_2$}
    \For{batch $\mathbf{b} \in \mathcal{D}_{\text{train}}$}
        \State Unfreeze all parameters
        \State $\boldsymbol{\alpha}, \mathbf{y}_{\text{aux}} \gets
        f_\theta(\mathbf{b})$ with AMP (FP16)
        \State $\mathcal{L} \gets$ Eq.~\eqref{eq:unified-loss}
        \State \textbf{if} $e > E_1 + 10$ \textbf{then}
            \State $\quad$ Enable gradient checkpointing (SIREN, OT aligner)
        \State $\theta \gets \theta - \eta \nabla_\theta\mathcal{L}$
    \EndFor
    \State \textbf{if} $\text{mIoU}_{\text{val}}$ plateaus for 10 epochs
        \textbf{then} $\eta \gets \eta \cdot 0.5$
\EndFor

\State \Return $\theta^* = \arg\max_\theta \text{mIoU}_{\text{val}}$
\end{algorithmic}
\end{algorithm}

\subsection{Implementation Statistics}

\subsection{Tensor Shape Trace}

Table~\ref{tab:shapes} provides the complete tensor shape transformation through
the V6 forward pass, tracing a single batch from input to output. This trace
is essential for verifying dimensional consistency across the 14 components.

\begin{table}[h]
\centering
\caption{Complete tensor shape trace through V6 forward pass}
\label{tab:shapes}
\begin{tabular}{lll}
\toprule
\textbf{Stage} & \textbf{Operation} & \textbf{Output Shape} \\
\midrule
Input & Optical & $(B, T, 10, H, W)$ \\
Input & SAR & $(B, T, 5, H, W)$ \\
Input & DEM & $(B, 5, H, W)$ \\
\midrule
1. Geometric & compute\_geometric\_invariants(DEM) & $(B, 5, H, W)$ \\
1. Geometric & GeometricInvariantEncoder(DEM) & $(B, 128, H, W)$ \\
2. Early Fusion & ModalNorm + Concat + Conv$1\times1$ & $(B, 128, H, W)$ \\
\midrule
3. Optical Enc. & Backbone(Opt) $\to$ FPN & $(B\times T, 512, H_4, W_4)$ \\
3. Optical FP2 & sp2 projection & $(B\times T, 256, H_2, W_2)$ \\
3. SAR Enc. s1 & SAR stem + stage1 + FiLM & $(B\times T, 64, H, W)$ \\
3. SAR Enc. s2 & stage2 + FiLM + downsample & $(B\times T, 128, H_2, W_2)$ \\
3. SAR Enc. s3 & stage3 + downsample & $(B\times T, 512, H_4, W_4)$ \\
\midrule
4. DEM$\to$Opt FiLM & DEMOpticalConditioner & $(B\times T, 512, H_4)$, $(B\times T, 256, H_2)$ \\
4. DEM$\to$Temp & MLP(AvgPool(dem)) & $(B\times H_4\times W_4, 512)$ \\
5. TemporalLite s1 & Conv1D$_{64}$ + gate & $(B\times H\times W, 64)$ \\
5. TemporalLite s2 & Conv1D$_{128}$ + gate & $(B\times H_2\times W_2, 128)$ \\
5. TemporalLite opt & Conv1D$_{256}$ + gate & $(B\times H_2\times W_2, 256)$ \\
5. Transformer & opt\_temporal + sar\_temporal & $(B, 512, H_4, W_4)$ each \\
\midrule
6. OT Align & JointOTGrassmannAligner & $(B, 512, H_4, W_4)$ each \\
7. CrossAttn H & CrossModalLite$_{64}$ & $(B, 64, H, W)$ \\
7. CrossAttn H/2 & CrossModalLite$_{128}$ & $(B, 128, H_2, W_2)$ \\
7. CrossAttn H/4 & CrossModalAttention$_{512}$ & $(B, 512, H_4, W_4)$ \\
8. Decoder & CARAFE$\uparrow\times2$ + SpatialRefine & $(B, 64, H, W)$ \\
\midrule
9. EDL Head & Conv$\to$GELU$\to$Conv + softplus & $(B, K, H, W)$ \\
11. LAI Head & Conv$\to$GAP$\to$MLP$\to$Softplus & $(B,)$ \\
11. Growth Head & Conv$\to$GAP$\to$MLP & $(B, 5)$ \\
11. Boundary Head & Conv$\to$Conv$1\times1$$\to$Sigmoid & $(B, 1, H, W)$ \\
11. Scene Head & Conv$\to$GAP$\to$Linear & $(B, 4), (B, K)$ \\
\bottomrule
\end{tabular}
\end{table}

Notation: $H_2 = H/2$, $H_4 = H/4$. All temporal modules first reshape
$(B, T, C, H, W) \to (B \times H \times W, T, C)$ (pixel-sequence format),
process temporally, then reshape back.

\subsection{Initialization Scheme}

All V6 components follow a unified initialization protocol derived from the
mathematical theory modules:

\begin{table}[h]
\centering
\caption{Per-component initialization scheme}
\begin{tabular}{lll}
\toprule
\textbf{Component} & \textbf{Method} & \textbf{Theory Basis} \\
\midrule
TemporalLite (Conv1D) & Spectral-Balanced: $\mathcal{N}(0,I) \to \|\mathbf{W}\|_2 = \sqrt{2}$ & Module 4, Thm. 3 \\
Optical Backbone & ImageNet/SeCo pre-trained + kaiming\_normal & Standard \\
SAR Encoder (stem) & kaiming\_normal(fan\_out) + BatchNorm $\gamma=1, \beta=0$ & Standard \\
SAR Encoder (FiLM) & $\gamma=0, \beta=0$ (identity init) & Module 1 \\
DEMOpticalConditioner & FiLM: $\gamma=0, \beta=0$ (identity) & Module 1 \\
CrossModalLite (QKV) & kaiming\_normal for depthwise + pointwise & Standard \\
CrossModalAttention & kaiming\_normal, GroupNorm $\gamma=1, \beta=0$ & Standard \\
EDL Head (Conv) & kaiming\_normal, final Conv $\mathcal{N}(0, 0.001)$ & EDL convention \\
Multi-Task Heads & kaiming\_normal, final Linear $\mathcal{N}(0, 0.01)$ & Standard \\
Transformer Encoder & trunc\_normal(std=0.02) + $\mathbf{W}_O = 0$ & GPT convention \\
DomainAdapter & $\gamma=1, \beta=0$ (identity mapping) & AdaIN convention \\
Placeholder params & trunc\_normal(std=0.02) & Standard \\
\bottomrule
\end{tabular}
\end{table}

\subsection{Regularization Summary}

\begin{table}[h]
\centering
\caption{Complete regularization configuration}
\begin{tabular}{lcc}
\toprule
\textbf{Technique} & \textbf{Value} & \textbf{Applied To} \\
\midrule
Dropout (EDL Head) & $p=0.3$ (Dropout2d) & Pre-classification features \\
Dropout (Temporal) & $p=0.1$ (timestep-wise) & Random time step masking \\
Dropout (Modality) & $p=0.1$ (per-sample) & Optical/SAR/DEM randomly dropped \\
Weight Decay & $\lambda=10^{-4}$ & All non-normalization parameters \\
Gradient Clipping & $\|\mathbf{g}\|_{\max}=1.0$ & All gradients \\
KL Annealing & $\lambda_{\max}=0.5$, linear 0$\to$50 epochs & EDL regularization \\
EMA (Temporal) & $\alpha=0.999$ & Consistency target projection \\
LayerNorm & $\epsilon=10^{-5}$ & TemporalLite, Transformer \\
GroupNorm & 32 groups, $\epsilon=10^{-5}$ & CrossModal, Decoder \\
BatchNorm & momentum=0.1 & SAR Encoder, DEM Encoder \\
\bottomrule
\end{tabular}
\end{table}

\subsection{Computational Complexity Analysis}

\begin{table}[h]
\centering
\caption{Per-component FLOPs and memory (input: $1\times12\times256\times256$)}
\begin{tabular}{lccc}
\toprule
\textbf{Component} & \textbf{FLOPs} & \textbf{Params} & \textbf{Peak Memory} \\
\midrule
Optical Backbone (ResNet50) & 4.1G & 25.6M & 1.2 GB \\
SAR Encoder & 1.8G & 2.1M & 0.4 GB \\
Geometric Invariants & 0.05G & 0.1M & 0.1 GB \\
TemporalLite $\times$ 3 & 0.02G & 3.7K & 0.05 GB \\
Transformer (opt + sar) & 2.4G & 8.4M & 1.8 GB \\
CrossModalAttention $\times$ 3 & 1.2G & 0.6M & 0.8 GB \\
OT Aligner (SGW + ADMM) & 0.3G & 0.3M & 0.3 GB \\
Decoder (CARAFE + SR) & 0.5G & 0.8M & 0.2 GB \\
Heads (EDL + 4 aux) & 0.1G & 0.6M & 0.1 GB \\
\midrule
\textbf{Total (FP32)} & \textbf{10.5G} & \textbf{49.0M} & \textbf{5.0 GB} \\
\textbf{Total (AMP FP16)} & \textbf{5.8G} & 49.0M & \textbf{2.8 GB} \\
\textbf{Total (+Grad Ckpt)} & 5.8G (compute) & 49.0M & \textbf{1.9 GB} \\
\bottomrule
\end{tabular}
\end{table}

With gradient checkpointing enabled on the Transformer encoder, OT aligner,
and SIREN DEM encoder, peak memory is reduced by $62\%$ (5.0 GB $\to$ 1.9 GB),
enabling training on consumer GPUs (RTX 3060 12GB, RTX 4070 12GB).

\subsection{3-Expert LateFusion Architecture}

The 3-Expert LateFusion (Block 9, P3) provides decision-level ensemble
fusion with vacuity-weighted voting:

\begin{equation}
    \mathbf{z}_{\text{final}} = w_{\text{opt}} \cdot \mathbf{z}_{\text{opt}} +
    w_{\text{sar}} \cdot \mathbf{z}_{\text{sar}} +
    w_{\text{fused}} \cdot \mathbf{z}_{\text{fused}}
\end{equation}

where the weights are computed from per-expert vacuity:
\begin{equation}
    u_m = \frac{K}{\sum_k \text{softplus}(z_m^{(k)}) + 1}, \quad
    \mathbf{w} = \text{softmax}(-\text{MLP}([u_{\text{opt}}, u_{\text{sar}}, u_{\text{fused}}]))
\end{equation}

Each expert is structurally identical (Conv$1\times1$$\to$GELU$\to$Conv$1\times1$,
256 hidden channels) but trained from independent weight initializations and
receives features from different modalities:

\begin{itemize}[nosep]
    \item \textbf{Expert\_opt:} receives optical-only features
    $\mathbf{f}_{\text{opt}} \in \mathbb{R}^{B \times 512 \times H_4 \times W_4}$
    \item \textbf{Expert\_sar:} receives SAR-only features
    $\mathbf{f}_{\text{sar}} \in \mathbb{R}^{B \times 512 \times H_4 \times W_4}$
    \item \textbf{Expert\_fused:} receives cross-modal fused features
    $\mathbf{f}_{\text{xm}} \in \mathbb{R}^{B \times 512 \times H_4 \times W_4}$
\end{itemize}

The key innovation over standard ensemble averaging is that vacuity weighting
automatically handles missing modalities: if optical data is unavailable,
$u_{\text{opt}} \to 1$ (maximum uncertainty), driving $w_{\text{opt}} \to 0$
via the softmax, and the ensemble gracefully degrades to SAR + fused experts.

\subsection{Self-Training Pipeline}

V6 employs iterative self-training with topological pseudo-label filtering:

\begin{algorithm}[h]
\caption{Topological Self-Training Loop}
\label{alg:self-training}
\begin{algorithmic}[1]
\Require Trained model $f_\theta$, unlabeled data $\mathcal{U}$,
         vacuity threshold $\tau_u^0 = 0.3$, confidence threshold $\tau_c^0 = 0.7$,
         max rounds $R = 5$
\Ensure Refined model $f_{\theta^*}$, pseudo-labeled dataset $\mathcal{P}$

\For{round $r = 1, \ldots, R$}
    \State $\mathcal{P}_r \gets \emptyset$
    \For{each unlabeled sample $\mathbf{x} \in \mathcal{U}$}
        \State $\boldsymbol{\alpha} \gets f_\theta(\mathbf{x})$
        \State $\mathbf{p} \gets \boldsymbol{\alpha} / \sum_k \alpha_k$
            \Comment{Class probabilities}
        \State $\mathbf{u} \gets K / \sum_k \alpha_k$
            \Comment{Scalar vacuity}

        \State \textbf{// Topological conflict classification (Module 3)}
        \State $t_{\text{conflict}} \gets \text{detect\_conflict\_type}(
            \boldsymbol{\alpha}_{\text{opt}}, \boldsymbol{\alpha}_{\text{sar}})$

        \State \textbf{// Three-way routing}
        \If{$t_{\text{conflict}} = \text{Noise}$ \textbf{and}
            $\max(\mathbf{p}) > \tau_c$ \textbf{and} $\mathbf{u} < \tau_u$}
            \State $\mathcal{P}_r \gets \mathcal{P}_r \cup \{(\mathbf{x},
                \arg\max(\mathbf{p}))\}$
            \Comment{Safe to pseudo-label}

        \ElsIf{$t_{\text{conflict}} = \text{Structural}$}
            \State $\boldsymbol{\alpha}' \gets \text{persistent\_correction}(
                \boldsymbol{\alpha})$
            \State \textbf{if} $\max(\boldsymbol{\alpha}' / \sum \boldsymbol{\alpha}') > \tau_c$
                \textbf{then}
                $\mathcal{P}_r \gets \mathcal{P}_r \cup \{(\mathbf{x},
                    \arg\max(\boldsymbol{\alpha}'))\}$
            \Comment{Corrected pseudo-label}

        \Else
            \State \textbf{skip}
            \Comment{HighOrder conflict — abstain}
        \EndIf
    \EndFor

    \State \textbf{// Fine-tune with pseudo-labels}
    \State $f_\theta \gets \text{train\_one\_epoch}(f_\theta,
        \mathcal{D}_{\text{train}} \cup \mathcal{P}_r)$

    \State \textbf{// Relax thresholds}
    \State $\tau_u \gets \min(0.5, \tau_u + 0.05)$
    \State $\tau_c \gets \max(0.7, \tau_c - 0.05)$
\EndFor

\State \Return $f_\theta, \bigcup_{r=1}^R \mathcal{P}_r$
\end{algorithmic}
\end{algorithm}

The three-way routing (Noise/Structural/HighOrder) replaces the traditional
binary vacuity threshold. Without topological conflict classification, the
baseline vacuity filter would accept \textbf{all} pixels with $u < \tau_u$,
including $17\%$ of hidden semantic contradictions ($\kappa < 0.1$ but
$H^1 \neq 0$) that degrade pseudo-label quality. The topological filter
explicitly rejects these, improving pseudo-label precision by $\sim$12\%
(Module 3 experimental data).

\subsection{Inference Procedure}

\subsubsection{Standard Inference}

During standard evaluation, V6 performs a single deterministic forward pass:
\begin{equation}
    \boldsymbol{\alpha} = f_\theta(\mathbf{X}_{\text{opt}}, \mathbf{X}_{\text{sar}},
    \mathbf{X}_{\text{dem}}, \mathbf{doy}), \quad
    \hat{y} = \arg\max_k \frac{\alpha_k}{\sum_j \alpha_j}
\end{equation}

\subsubsection{Monte Carlo Dropout Uncertainty}

For uncertainty quantification, $N$ stochastic forward passes are performed
with Dropout2d enabled in the EDL head:
\begin{equation}
    \boldsymbol{\alpha}^{(n)} = f_\theta(\mathbf{X}; \text{dropout}=0.3),
    \quad n = 1, \ldots, N
\end{equation}
Fused evidence: $\boldsymbol{\alpha}_{\text{fused}} = \frac{1}{N}\sum_n
\boldsymbol{\alpha}^{(n)}$.

Uncertainty decomposition (from \texttt{dirichlet\_to\_predictions}):
\begin{align}
    \text{vacuity}(\mathbf{x}) &= \frac{K}{\sum_k \alpha_{\text{fused},k}}
        \quad \text{(evidence insufficiency)} \\
    \text{dissonance}(\mathbf{x}) &= 1 - \sum_k \left(\frac{\alpha_k}{S}\right)^2
        \quad \text{(class conflict)}
\end{align}

\subsubsection{Test-Time Augmentation (TTA)}

Optional horizontal flip TTA:
\begin{equation}
    \boldsymbol{\alpha}_{\text{TTA}} = \frac{1}{2}\left(
    \boldsymbol{\alpha}(\mathbf{X}) + \text{flip}(\boldsymbol{\alpha}(
    \text{flip}(\mathbf{X})))\right)
\end{equation}

\subsubsection{Missing Modality Handling}

V6 is robust to missing modalities through placeholder parameters:
\begin{align}
    \tilde{\mathbf{X}}_{\text{opt}} &=
    \begin{cases}
        \mathbf{X}_{\text{opt}} & \text{if optical available} \\
        \mathbf{p}_{\text{opt}} \sim \mathcal{N}(0, 0.02) & \text{otherwise}
    \end{cases}
\end{align}
where $\mathbf{p}_{\text{opt}} \in \mathbb{R}^{1 \times 512 \times 1 \times 1}$
is a learned placeholder expanded to match spatial dimensions. During training,
modality dropout ($p=0.1$) randomly drops one modality to learn robust
placeholders, controlled by \texttt{modality\_mask} parameter.

The 3-Expert LateFusion automatically down-weights experts whose input modality
is missing (vacuity $\to 1 \to$ weight $\to 0$), providing graceful degradation
without explicit modality-availability flags at inference time.

\subsection{Implementation Statistics}

\begin{table}[h]
\centering
\caption{V6 model statistics}
\begin{tabular}{lcc}
\toprule
\textbf{Component} & \textbf{Parameters} & \textbf{Source} \\
\midrule
Backbone (ResNet50) & $\sim$25M & timm \\
TemporalLite $\times$ 3 & $\sim$3.7K & temporal\_lite.py \\
Multi-Scale CrossAttn & $\sim$0.3M & \_base.py \\
DEM New Paths (3) & $\sim$0.25M & \_base.py \\
Multi-Task Heads (3) & $\sim$15K & multi\_task\_heads.py \\
LightSceneHead & $\sim$18K & \_base.py \\
EDL Head & $\sim$0.5M & fusion\_net\_v5\_edl.py \\
\textbf{Total} & $\mathbf{\sim49M}$ & $+4\text{M}$ (9\%) over V5Pro \\
\bottomrule
\end{tabular}
\end{table}

The complete codebase consists of 22 model files totaling 7,329 lines, with
102 classes, 322 functions, and 313 passing tests. All mathematical modules
are integrated as drop-in components with opt-in activation.
